\documentclass[a4paper,11pt]{article}

\usepackage[utf8]{inputenc}
\usepackage[T1]{fontenc}
\usepackage[french]{babel}

\usepackage[top=25mm,left=40mm,right=40mm]{geometry}
\usepackage{fourier}
\usepackage[babel]{csquotes}
\usepackage[hyphens]{url}
\usepackage{hyperref}
\hypersetup{
    colorlinks,
    citecolor=black,
    filecolor=black,
    linkcolor=black,
    urlcolor=black
}
\usepackage[normalem]{ulem}

\title{TP3 : Clustering et word2vec}
\author{Traitement automatique de corpus}
\date{4 novembre 2025}

\begin{document}
\maketitle
\thispagestyle{empty}

Sur la base des éléments méthodologiques et des enseignements techniques présentés lors du cours théorique, il est demandé dans le cadre de ce TP :

\begin{itemize}
\item d'effectuer un clustering des documents CAMille pour une décennie au choix et d'interpréter les résultats
\item d'entraîner un modèle word2vec sur l'ensemble du corpus et d'explorer les relations entre vecteurs
\end{itemize}

Pour ce faire, vous utiliserez différentes librairies Python vues au cours comme scikit-learn et gensim.

\vspace{0.7cm}

Les étapes à mettre en \oe uvre sont les suivantes :

\begin{enumerate}
\item Choisissez une décennie (1890--1899, 1900--1909, 1910--1919, etc.)
\item Effectuez un clustering des documents de cette décennie grâce au notebook \texttt{s2\_clustering.ipynb}, en adaptant éventuellement le nombre de clusters désirés
\item Tentez d'interpréter les résultats obtenus : les clusters semblent-ils faire sens ? Si oui/non, comment l'expliquez-vous ? Aidez-vous au besoin d'une analyse de keywords et/ou wordcloud (voir TP2).
\item Téléchargez sur l'UV le fichier zippé \texttt{sents.txt}, déjà segmenté en phrases, et placez-le dans le dossier data
\item Entraînez un modèle word2vec (word embeddings) sur ces phrases grâce au notebook \texttt{s3\_word\_embeddings.ipynb}, en adaptant éventuellement les paramètres window (taille de la fenêtre) et min\_count (nombre minimum d'occurrences d'un mot)
\item Vous pouvez entraîner plusieurs modèles afin de comparer leurs performances, en procédant par essais-erreurs pour choisir le meilleur modèle
\item Explorez le modèle retenu à l'aide des fonctions \texttt{similarity} et \texttt{most\_similar} (choisissez au moins trois exemples pour chaque fonction)
\item Rassemblez votre code dans un nouveau notebook \texttt{tp3.ipynb} placé dans le dossier \enquote{tp3} (n'oubliez pas de pusher)
\item Soumettez sur l'UV un court rapport (3-4 pages, format PDF) présentant votre méthodologie et les résultats obtenus
\end{enumerate}

Pour toute question, n'hésitez pas à me contacter via \texttt{max.de.wilde@ulb.be} en mettant toujours Denis Lebailly en copie (\texttt{denis.lebailly@ulb.be}).

\end{document}